\section{Requirement}


\begin{center}
	\begin{tikzpicture}
		% Hộp (rectangle)
		\node[
		draw=violet!70!black,
		fill=violet!10,
		rounded corners=2pt,
		inner xsep=6pt,
		inner ysep=6pt,
		text width=0.85\textwidth,
		align=justify
		] (box) {
			\small
			Thiết kế một bộ \textbf{Floating Point (FPU)} có thể cộng và trừ hai số 
			\textbf{floating-point 32-bit}. Phương pháp thiết kế sử dụng thuần là bộ tổ hợp.
		};
	\end{tikzpicture}
\end{center}

\noindent Sau khi thiết kế thành công, cần:

\begin{itemize}
	\item Viết chương trình kiểm định thiết kế \textbf{FPU} trên.
	\item Thực hiện thiết kế trên KIT FPGA để kiểm tra chức năng của thiết kế.
	\item Tiến hành tổng hợp thiết kế sử dụng các standard cell trong thư viện có sẵn tại các trường hợp thực tế và loại bỏ các glitchs.
\end{itemize}

\noindent Ngõ vào và ngõ ra của thiết kế được mô tả như sau:

\begin{table}[H] % or [t], [b], [!ht], etc.
	\centering
	\begin{tabular}{|>{\centering\arraybackslash}m{2cm}
		                 |>{\centering\arraybackslash}m{2cm}
		                 |>{\centering\arraybackslash}m{2cm}
		                 |>{\raggedright\arraybackslash}m{7.5cm}|}
	     \hline
	     Name 			& Port Type 	& Width 		& Description\\
	     \hline
	     i\_32\_a		& Input 		& 32			& 32-bit floating point number A\\
	     \hline
	     i\_32\_b		& Input 		& 32			& 32-bit floating point number B\\
	     \hline
	     i\_add\_subb	& Input 		& 1				& Operation: 0 = add, 1 = sub\\
	     \hline
		 o\_32\_s		& Output 		& 32			& 32-bit floating point output\\
	     \hline
	     o\_ov\_flag	& Output 		& 1				& Overflow Flag: 1 = overflow, 0 = no overflow occurs\\
	     \hline
	     o\_un\_flaag 	& Output 		& 1 			& Underflow Flag: 1 = underflow, 0 = no underflow occurs\\
	     \hline
	 \end{tabular}
	 \caption{Input and Output Ports.}
	 \label{tab: input-and-output-ports}
\end{table}